% THEORY SECTION - Phase 2 Implementation
% Approximately 800 words
% Theoretical Foundations for Sentiment-Based Forecasting

\section{Theoretical Foundations of Narrative and Sentiment in Economics}

Why do narratives move markets and macroeconomies? This section develops the theoretical case for sentiment-based forecasting, tracing mechanisms from behavioral psychology through expectation formation to macroeconomic outcomes. Understanding these channels answers why sentiment forecasts should work and clarifies when they might mislead.

\subsection{Why Do Narratives Matter? Channels of Influence}

Economic agents form beliefs under radical uncertainty. Future GDP, inflation, and asset prices are unknowable; official statistics lag months behind events; economic models conflict. In this epistemic void, narratives provide coherent explanations and expectations for otherwise opaque economic futures.

Behavioral finance emphasizes how narratives drive animal spirits---waves of investor confidence and pessimism disconnected from fundamentals. Shiller (2017) argues that successful narratives are those that ``go viral,'' spreading through social networks and media, shaping expectations through narrative plausibility rather than empirical proof. The dot-com bubble, housing bubble, and cryptocurrency manias exemplify narratives divorcing prices from fundamental value for extended periods. Sentiment indices capturing these narratives should forecast reversals when reality eventually confronts fantasy.

Macroeconomic expectations channels provide a complementary mechanism. Modern Phillips curves and consumption functions emphasize expectations: inflation expectations drive wage-setting and pricing behavior (Phillips curve), consumer confidence drives spending plans (consumption Euler equations). Through these channels, narratives shaping expectations directly influence real outcomes. When media turns pessimistic, consumer confidence falls, dampening demand and inflation pressure; firms cut investment based on gloomy revenue forecasts. Angeletos and La'O (2013) formalize this: in their general equilibrium model, sentiment shocks (beliefs about future productivity and demand) drive business cycles even without fundamental shocks. The mechanism is expectation spillovers: if agents believe others are pessimistic, they become pessimistic, self-fulfilling the narrative.

Information aggregation provides a third channel. Markets aggregate dispersed private information into prices; media aggregates newsworthy events and expert interpretation. If media narratives reflect underlying economic realities that official statistics miss, sentiment indices serve as real-time information aggregators. Tetlock (2007) proposes this mechanism: media pessimism predicts stock returns because media pessimism reflects investor caution before fundamental downturns become official. Under this interpretation, sentiment is merely a summary statistic for information already in prices; its predictive power reflects information's lag into official statistics.

These three channels---animal spirits (sentiment overreaction), expectations transmission, and information aggregation---are empirically indistinguishable in many contexts. A forecaster observing strong sentiment predictive power cannot determine which mechanism dominates: Is it behavioral overreaction likely to reverse, or is it genuine information predictive of future fundamentals? This ambiguity matters for practice: overreaction suggests contrarian strategies; information aggregation suggests trend-following strategies. The literature provides evidence for both mechanisms in different contexts and time periods.

\subsection{Sentiment Versus Noise: The Fundamental Debate}

The core empirical debate in the literature concerns sentiment's fundamental nature. Are sentiment-driven asset price movements driven by fundamentals, or are they noise that will eventually reverse?

Tetlock (2007), examining Wall Street Journal columns, finds that high media pessimism predicts negative stock returns 1--5 trading days ahead, then reversals: pessimistic columns followed by positive returns 15--40 days later. This V-shaped return pattern suggests overreaction: sentiment shocks initially move prices away from fundamental value, then market corrections restore equilibrium. The predictive power reflects not information about future earnings, but temporary sentiment-driven mispricings.

In contrast, Da, Engelberg, and Gao (2015) construct the FEARS index from internet search volume for fear-related terms. They find that FEARS negatively predicts returns one day ahead (pessimism coincides with selling) but positively predicts returns 2--3 weeks ahead (delayed recovery). The reversal pattern parallels Tetlock, yet the economic interpretation differs: if agents search for fear-related terms when fundamentals deteriorate, then high FEARS indicates low asset values, and subsequent positive returns reflect price recovery as bad news prices in. Under this interpretation, sentiment is information, not noise.

These competing interpretations persist because both can fit return patterns: noise reverses in 2--4 weeks; information takes 1--2 weeks to incorporate into prices. Without exogenous variation (IV strategies, natural experiments) severing sentiment from fundamentals, causality remains ambiguous.

Contextual variation offers partial resolution. During crisis periods (2008--2009, COVID-19 March 2020), sentiment predictive power surges, suggesting heightened information content when uncertainty is extreme. During normal periods, sentiment effects are modest and reversals frequent, suggesting noise dominance. The implication: sentiment forecasting is regime-dependent; practitioners should monitor sentiment effectiveness within their forecast period and adjust confidence accordingly.

\subsection{The Causality Question: Prediction Versus Causation}

The vast majority of studies (55/66 papers in the review) demonstrate that sentiment \emph{predicts} economic outcomes: after controlling for lagged fundamentals, sentiment Granger-causes inflation, GDP, or stock returns. Yet prediction does not imply causation in policy or scientific contexts. Understanding whether sentiment \emph{causes} outcomes matters for: (a) interpretation of mechanisms, (b) policy implications (can governments move sentiment to influence real outcomes?), and (c) equilibrium modeling (sentiment shocks should enter structural macro models).

Establishing causation requires either: (1) natural experiments isolating exogenous sentiment variation, or (2) instrumental variables severing sentiment from confounders. The literature provides few examples. Hansen and McMahon (2016) exploit timing of FOMC statements as quasi-exogenous shocks to policy sentiment. They find that shocks to forward guidance (sentiment about future policy) cause reactions in financial markets within hours of publication, supporting causality. However, they do not establish whether sentiment-driven financial reactions causally influence real economic outcomes (GDP, inflation); only that sentiment causally influences market sentiment, a weaker claim.

Causal identification remains largely absent from economic sentiment literature. Most studies employ Granger causality, temporal precedence, or prediction-in-real-time, all compatible with reverse causality or confounding. For example: deteriorating GDP forecasts could drive media pessimism (GDP causes sentiment), not vice versa. Similarly, financial wealth shocks (stock market decline) could drive both pessimism and economic outcomes (wealth effect), making causality appear to run from sentiment to outcomes when both are joint consequences of wealth shocks.

The implication: sentiment-based forecasting succeeds empirically because sentiment contains information about future outcomes, but the literature does not establish whether that information reflects sentiment's causal influence or mere correlation with contemporaneous unobserved drivers of outcomes. For practitioners, this ambiguity is acceptable: forecasting only requires conditional prediction, not causality. For policymakers, it is problematic: cannot infer whether sentiment management (central bank communication, policy guidance) causally influences real outcomes without exogenous identification.

\subsection{Conceptual Framework: Narratives in Economic Forecasting}

Integrating the above elements, we propose a framework for understanding sentiment-based forecasting (Figure~\ref{fig:framework} in main text):

\noindent
Narratives $\rightarrow$ Expectations $\rightarrow$ Behavior $\rightarrow$ Outcomes

Narratives (media stories, analyst commentary, policy communication) shape agent expectations through plausibility and resonance. Expectations influence behavior: consumption, investment, wage demands, prices. Behavior aggregates into macroeconomic outcomes: GDP, inflation, asset prices. Sentiment indices measure narratives at scale (``how pessimistic is media today?''), providing a real-time proxy for expectation dynamics.

This framework clarifies when sentiment forecasting should work: when narratives are the binding constraint on expectation formation (high uncertainty environments, new crises, breakdown of precedent), sentiment is highly informative; when agents update on hard data (low unemployment reducing unemployment fear), sentiment lags and has little predictive power. The framework also clarifies limitations: sentiment indices measure narratives in covered sources (media, surveys), missing narratives in private conversations, messaging platforms, or face-to-face communication that may be equally economically important.

The Narrative Economics program (Shiller, 2017; Benigsen et al., 2022) operationalizes this framework by studying which narratives spread, how they interact with economic outcomes, and how they drive economic history. Recent advances in sentiment extraction make this research program computationally tractable: instead of qualitative narrative analysis, researchers deploy neural networks to identify narrative themes across millions of documents, enabling large-scale tests of narrative-outcome connections.

